\documentclass[11pt,a4paper]{ctexart}
\usepackage[margin=2.45cm]{geometry}
\usepackage{amsmath,amssymb,amsthm,mathtools}
\usepackage{booktabs}
\usepackage{float}          % 【补】提供 [H] 精确位置（table、algorithm 均需）
\usepackage{algorithm}      % 【补】algorithm 浮动环境
\usepackage{algpseudocode}  % 【补】\State/\If/\For/\While/\Comment/\Call/\Statex/\Return
\usepackage{xcolor}
\usepackage{listings}       % 【补】lstlisting 代码环境
\usepackage[hidelinks]{hyperref}   % hyperref 尽量放在最后
\usepackage{microtype}
% \usepackage{verbatim}    % 原文加载但全文未使用，已删

% 【补】正文多处使用 \Ord_m(2)（2 模 m 的阶），原文未定义，必须声明
\DeclareMathOperator{\Ord}{Ord}

% 【可选】浮动体与代码清单标题中文化
\floatname{algorithm}{算法}
\renewcommand{\lstlistingname}{代码清单}

% 【补】代码清单样式：支持中文注释与长行自动折行（建议 XeLaTeX 编译）
\lstset{
  language=Python,
  basicstyle=\ttfamily\footnotesize,
  keywordstyle=\color{blue!65!black},
  commentstyle=\color{gray},
  stringstyle=\color{teal!75!black},
  breaklines=true,
  columns=flexible,
  keepspaces=true,
  showstringspaces=false,
  frame=single,
  framerule=0.3pt,
  tabsize=4
}

\newtheorem{theorem}{定理}
\newtheorem{lemma}[theorem]{引理}
\newtheorem{proposition}[theorem]{命题}
\newtheorem{definition}[theorem]{定义}
\theoremstyle{remark}
\newtheorem{remark}[theorem]{注}
\newcommand{\Q}{\mathbb Q}
\newcommand{\Z}{\mathbb Z}
\newcommand{\Norm}{\operatorname N}   % 未使用，可删
\newcommand{\vp}{v_{\mathfrak p}}     % 未使用，可删

\title{方程 \(2^n-n=k^2\) 的全部正整数解只有 \(n=1,7\)}
\author{}
\date{2026年9月}

\begin{document}
\maketitle

\begin{abstract}
本文把基于 Ridout 定理的非有效有限性证明替换为一个有效论证。
通过在 \(\Z[\sqrt2]\) 中用基本单位 \(3+2\sqrt2\) 约化，并应用 Matveev 的显式对数线性形式下界，证明每个正整数解均满足
\[ n\le 4\,850\,864\,607\,069\,127. \]
再依据平方数模 $m$ 的值在数值上界内筛选，可行解仅有 $n=1,7$.
\end{abstract}

\section{问题与初等限制}
考虑
\begin{equation}\label{eq:main}
2^n-n=k^2,\quad n,k\in\Z_{>0}.
\end{equation}
若 \(n\ge4\) 为偶数，令 \(x=2^{n/2}\)，则
\[ (x-k)(x+k)=n. \]
两因子均为正整数，故 \(x+k\le n\)；但 \(x\ge n\)，矛盾。
直接检查 \(n=2\) 也不成立。因此除解 \((n,k)=(1,1)\) 外，可写
\begin{equation}\label{eq:mq}
n=2m+1,\quad m\ge1,\quad q=2^m.
\end{equation}
原方程化为
\begin{equation}\label{eq:norm}
k^2-2q^2=-n.
\end{equation}
还有一个便于后续筛选的必要条件：奇数平方模 \(8\) 等于 \(1\)，
故 \(n\ge3\) 时
\begin{equation}\label{eq:mod8}
n\equiv7\pmod 8.
\end{equation}

\section{在 \texorpdfstring{\(\Z[\sqrt2]\)}{Z[sqrt(2)]} 中约化}
记
\[ A=k+q\sqrt2,\quad A'=k-q\sqrt2,\quad \eta=3+2\sqrt2. \]
\begin{remark}
方程的解必须满足 $k=\lfloor2^{\frac{n}{2}}\rfloor=\lfloor q\sqrt{2}\rfloor$，
因此 $-A'$ 为 $2^{\frac{n}{2}}$ 的小数部分，
我最开始的思路是限制它的上界 $O(n)$ 和下界 $\Theta(n)$，在 $n$ 充分大时导出矛盾，或许可以让本证明有一些来由。
\end{remark}

由 \eqref{eq:norm}，有 \(AA'=-n\)。这里 \(A>0>A'\)，而 \(\eta\) 是 \(\Q(\sqrt2)\) 中范数为 \(1\) 的基本正单位。
由反证法可知，存在唯一整数 \(t\ge0\)，使
\begin{equation}\label{eq:reduction}
\beta=A\eta^{-t},\quad \sqrt n\le\beta<\eta\sqrt n.
\end{equation}
因为 \(\eta^{-1}=3-2\sqrt2\in\Z[\sqrt2]\)，故 \(\beta\in\Z[\sqrt2]\)。其共轭满足
\begin{equation}\label{eq:beta-conj}
\beta'=-\frac n\beta,\quad |\beta'|\le\sqrt n.
\end{equation}
此外 \(t<m\)。事实上，
\[ \eta^t=\frac A\beta\le \frac{A}{\sqrt{n}}< \frac{2^{m+3/2}}{\sqrt{2m+1}}<\eta^m. \]
最后一个不等式在 \(m=1\) 时可直接检查；在 \(m\ge2\) 时由 \(\eta>4\) 立即得到。因此下文 Matveev 定理中的系数上界可以取 \(B=m\)。

\section{实对数线性形式}
定义
\begin{equation}\label{eq:Delta}
\Delta=\eta^t2^{-m}\frac{\beta}{2\sqrt2}-1
=\frac{A}{2q\sqrt2}-1.
\end{equation}
由于 \(k<q\sqrt2\)，有 \(-1<\Delta<0\)，特别地 \(\Delta\ne0\)。
再由 \(A'=-n/A\) 以及 \(A>q\sqrt2\)，得到
\begin{equation}\label{eq:upperDelta}
0<|\Delta|=\frac{|A'|}{2q\sqrt2}
=\frac{n}{2q\sqrt2 A}
<\frac{n}{4q^2}=\frac{n}{2^{2m+2}}.
\end{equation}

\begin{definition}[绝对对数高度 \cite{NesterenkoWaldschmidt}]\label{def:absLogHeight}
设 $\alpha\neq 0$ 为代数数，其在 $\mathbb{Q}$ 上的次数为 $d$，最小多项式为
\[ f(X)=a_0\prod_{i=1}^{d}\bigl(X-\alpha^{(i)}\bigr)\in\mathbb{Z}[X], \]
其中 $\alpha^{(i)}$（$1\le i\le d$）为 $\alpha$ 的全部复共轭。则 $\alpha$ 的\emph{绝对对数高度}定义为
\begin{equation}\label{eq:absLogHeight}
h(\alpha)=\frac{1}{d}\Bigl(\log|a_0|+\sum_{i=1}^{d}\log\max\{1,\,|\alpha^{(i)}|\}\Bigr).
\end{equation}
\end{definition}

由 \eqref{eq:reduction}--\eqref{eq:beta-conj} 以及 \(\beta\) 为二次代数整数，
\begin{equation}\label{eq:hbeta}
h(\beta)
=\frac12\left(\log\max\{1,\beta\}
+\log\max\{1,|\beta'|\}\right)
\le\frac{\log n+\log\eta}{2}.
\end{equation}
令 \(\gamma=\beta/(2\sqrt2)\)。高度的次可加性 $h(xy)\le h(x)+h(y)$ 给出
\begin{equation}\label{eq:hgamma}
2h(\gamma)
\le \log n+\log\eta+2\log(2\sqrt2).
\end{equation}
记
\begin{equation}\label{eq:c0}
c_0=\log\eta+2\log(2\sqrt2)
=3.842188\ldots<3.843.
\end{equation}

我们采用如下版本的 Matveev 显式估计，在参考文献 \cite{Matveev} 的 \textbf{Corollary 2.3} 中：若 \(D\) 次数域中的正实
代数数 \(\alpha_1,\ldots,\alpha_s\) 与整数 \(b_i\) 使
\(\Lambda=b_1\ln{\alpha_1}+\cdots+b_s\ln{\alpha_s}\ne0\)，并且
\[ A_i\ge\max\{Dh(\alpha_i),|\log\alpha_i|,0.16\}, \quad
B\ge\max\{|b_1|,\ldots,|b_s|\}, \]
则
\begin{align}\label{eq:matveev}
\log|\Lambda| {}>&-C_1(s)30^{s+3}s^{4.5}D^2(1+\log D)(1+\log B) A_1\cdots A_s.
\end{align}
其中
\[ C_1(s)=\min\{0.5e\,30^{s+3}s^{4.5}, 2^{6s+20}\}. \]

在 Matveev 估计中取
\[ s=3,\quad D=2,\quad
(\alpha_1,\alpha_2,\alpha_3)=(\eta,2,\gamma),\quad (b_1,b_2,b_3)=(t,-m,1), \]
并取
\[ A_1=\log\eta,\quad A_2=2\log2,\quad A_3=\log n+c_0,\quad B=m. \]
由 \eqref{eq:hgamma} 可知这些选择合法。于是
\begin{equation}\label{eq:lowerDelta}
\log|\Lambda|> -F(1+\log m)(\log n+c_0),
\end{equation}
其中
\begin{align*}
F={}&0.5e\,30^6 3^{4.5}\,4(1+\log2) (\log\eta)(2\log2)\\
={}&2.300582689682\ldots\times10^{12}
<2.31\times10^{12}.
\end{align*}

\begin{lemma}\label{lem:log-ineq}
设 $x\in(-1, 0)$，则
\[ \ln(x+1)>\frac{x}{x+1}. \]
\end{lemma}

\begin{proof}
构造辅助函数
\[ f(x)=\ln(x+1)-\frac{x}{x+1},\qquad x\in(-1,0), \]
则只需证 $f(x)>0$。对 $f$ 求导：
\[ f'(x)=\frac{1}{x+1}-\frac{1}{(x+1)^2}=\frac{x}{(x+1)^2}. \]
当 $x\in(-1,0)$ 时，分子 $x<0$，分母 $(x+1)^2>0$，故 $f'(x)<0$，
即 $f$ 在 $(-1,0)$ 上严格单调递减。
由于 $f$ 在 $x=0$ 处连续，且 $f(0)=\ln 1-0=0$。任取 $x\in(-1,0)$，
由拉格朗日中值定理，存在 $\xi\in(x,0)$ 使得
\[ f(0)-f(x)=f'(\xi)\,(0-x). \]
因 $\xi\in(-1,0)$，有 $f'(\xi)<0$；又 $0-x>0$，故
\[ f(0)-f(x)<0 \quad\Longrightarrow\quad f(x)>f(0)=0, \]
即 $\ln(x+1)>\dfrac{x}{x+1}$。引理得证。
\end{proof}

已知 \(\Delta=\exp(\Lambda)-1\) 且 \(|\Delta|<1/2\)，在引理~\ref{lem:log-ineq} 中令 $x=\Delta$，可得到
\[ |\Lambda|=-\Lambda<-\frac{\Delta}{\Delta+1}=\frac{|\Delta|}{1-|\Delta|}<2\,|\Delta| \]
因此
\[ \log|\Delta| > \log\left(\frac{1}{2}|\Lambda|\right) = \log|\Lambda| - \log 2. \]
结合 \eqref{eq:upperDelta} 和 \eqref{eq:lowerDelta}，每个解必须满足
\begin{equation}\label{eq:keyineq}
(2m+1)\log2-\log(2m+1)
<2.31\times10^{12}(1+\log m)
\bigl(\log(2m+1)+3.843\bigr).
\end{equation}

\section{数值截断}
令 \(M=2.425432303534564\times10^{15}\)，并对实数 \(x\ge M\) 定义
\begin{align*}
G(x)={}&(2x+2)\log2-\log(2x+1)\\
&-2.31\times10^{12}(1+\log x)
\bigl(\log(2x+1)+3.843\bigr).
\end{align*}
直接代入（所有舍入均朝有利于下界的方向）得到
\[ G(M)>0.30>0. \]
同时
\begin{align*}
G'(x)={}&2\log2-\frac2{2x+1}\\
&-2.31\times10^{12}\left(
\frac{\log(2x+1)+3.843}{x}
+\frac{2(1+\log x)}{2x+1}\right).
\end{align*}
在 $x\geq M$ 上是增函数，并且 $G'(M)\approx 1.31 >0$，所以 \(G(x)>0\) 对全部 \(x\ge M\) 成立，
这与 \eqref{eq:keyineq} 矛盾。

\begin{theorem}\label{thm:bound}
方程 \eqref{eq:main} 的每个正整数解均满足
\[ \boxed{n\le 4\,850\,864\,607\,069\,127}. \]
\end{theorem}

\begin{proof}
解 \(n=1\) 显然满足结论。其余解有 \(n=2m+1\)、\(m\ge1\)。
上述论证说明 \(m<M=2.425432303534564\times10^{15}\)。由于 \(m\) 为整数，
\(n=2m+1\le 4\,850\,864\,607\,069\,127\)，结论随即得到。
\end{proof}

\section{筛除不可行解}
本节将“$2^n - n$ 是否为完全平方”的判定从任意位长整数的高精度开方，改造为 $O(1)$ 的查表运算：
以初等数论的必要性条件排除绝大多数 $n$（每一次排除均自带
可独立复核的否定证书），仅对极少数幸存者执行一次精确判定。流程的正确性不依赖
任何概率假设；唯一的非认证环节是对幸存者的最终精确复核
（算法~\ref{alg:sieve} 第~8--11 行）。

\subsection{筛的数学原理}

\paragraph{必要性条件。}
对正整数 $m$，定义模 $m$ 的平方剩余集
\begin{equation}\label{eq:qr-def}
QR(m) \coloneqq \{\, x^2 \bmod m : x \in \mathbb{Z} \,\},
\end{equation}
其中 $0=0^2$ 必须计入。若 $N = k^2$ 为完全平方，则对一切 $m$
\begin{equation}\label{eq:necessity}
N \bmod m = (k \bmod m)^2 \bmod m \;\in\; QR(m).
\end{equation}
其逆否命题——只要存在一个 $m$ 使 $N \bmod m \notin QR(m)$，则 $N$ 非平方——
即是筛的合法性依据。

\paragraph{一个简单的必要性条件。}
\begin{equation}\label{eq:mod8.1}n \equiv 7 \pmod 8\end{equation}

\paragraph{谓词的有限周期性。}
被检验量 $2^n - n$ 中 $n$ 出现两次——指数上的与被减的——分别服从不同周期。记
\begin{equation}\label{eq:ord}
\Ord_m(2) \coloneqq \min\{\, d \ge 1 : 2^d \equiv 1 \pmod m \,\},
\end{equation}
则幂部分满足
\begin{equation}\label{eq:period}
2^{n} \equiv 2^{\,n \bmod \Ord_m(2)} \pmod m,
\end{equation}
而减数部分只依赖 $n \bmod m$。对奇素数 $p$，由 Lagrange 定理 $\Ord_p(2) \mid p-1 < p$，
故 $\gcd\bigl(p,\Ord_p(2)\bigr)=1$，中国剩余定理给出
\begin{equation}\label{eq:crt-L}
L \coloneqq \operatorname{lcm}\bigl(p,\ \Ord_p(2)\bigr) = p \cdot \Ord_p(2),
\qquad
\mathbb{Z}_L \;\cong\; \mathbb{Z}_p \times \mathbb{Z}_{\Ord_p(2)}.
\end{equation}
于是“是否为模 $p$ 剩余”这一谓词坍缩为 $\mathbb{Z}_L$ 上的布尔表
\begin{equation}\label{eq:table-def}
G_p[t] \coloneqq
\Bigl[\, \bigl(2^{t} - t\bigr) \bmod p \in QR(p) \Bigr] \in \{0,1\},
\qquad t \in \mathbb{Z}_L,
\end{equation}
且对一切 $n \ge 0$
\begin{equation}\label{eq:equiv}
\underbrace{\Bigl[\,(2^n - n)\bmod p \in QR(p)\Bigr]}_{\text{原谓词，依赖整个 } n}
\;=\; G_p\bigl[\,n \bmod L\,\bigr].
\end{equation}
式~\eqref{eq:equiv} 是筛能处理 $n \sim 10^{15}$ 的根源：单点检验从 $O(\log n)$ 次模乘降为一次取模加一次查表，成本与 $n$ 的位数无关。
（合数模，如轮子中的 $m=9$：此时 $\gcd\bigl(m,\Ord_m(2)\bigr)$ 可大于 $1$，
周期取 $L=\operatorname{lcm}\bigl(m,\Ord_m(2)\bigr)$；又因 $QR(9)$ 模 $3$ 的投影落在 $QR(3)$ 内，“筛 9”蕴含“筛 3”，故轮子取 $9$ 而非 $3$。）

\paragraph{单筛淘汰率。}
对奇素数 $p$，固定 $r \in \mathbb{Z}_{\Ord_p(2)}$：集合
$\{t \in \mathbb{Z}_L : t \equiv r \ (\mathrm{mod}\ \Ord_p(2))\}$ 的 $p$ 个元素
在模 $p$ 意义下两两不同，而 $s \mapsto 2^r - s$ 是 $\mathbb{Z}_p$ 上的双射，
故该集合的判值取遍 $\mathbb{Z}_p$，其中恰 $|QR(p)| = (p+1)/2$ 个剩余。于是
\begin{equation}\label{eq:count}
\frac{\#\{\,t \in \mathbb{Z}_L : G_p[t] = 1\,\}}{L}
\;=\; \frac{p+1}{2p} \;=\; \frac12 + \frac{1}{2p},
\end{equation}
即每个素数\emph{精确}淘汰 $\dfrac{p-1}{2p}$ 的候选——是定理而非启发式。

\paragraph{否定证书。}
若 $n$ 被素数 $p$ 筛除，记 $a = (2^n - n) \bmod p$，由 Euler 判据
\begin{equation}\label{eq:euler}
a^{(p-1)/2} \equiv \Bigl(\frac{a}{p}\Bigr) \in \{+1,-1\} \pmod p,
\qquad -1 \iff a \notin QR(p),
\end{equation}
任何第三方都可用 $O(\log n)$ 的快速幂独立复核“$a$ 非剩余 $\Rightarrow$ $2^n - n$ 非平方”。
证书体积 $O(\log p)$ 字节，与 $n$ 无关；对比之下
“肯定证书”（平方根本身）约 $n/2$ 比特，$n \sim 10^{15}$ 时约 62\,TB——
这解释了本筛为何只承诺廉价地证明\emph{不是}平方。

\paragraph{无假阴性。}
由式~\eqref{eq:necessity}，若 $2^n - n = k^2$，则对\emph{每一个}筛
（素数或合数）都有 $G[\,n \bmod L\,] = 1$：真解永远存活。
因此筛只可能产生假阳性（由最终精确复核消除），不存在假阴性。

\subsection{算法}
算法分三段：单模数建表（算法~\ref{alg:table}）、轮子构建
（算法~\ref{alg:wheel}）、主筛与精确复核（算法~\ref{alg:sieve}）。

\begin{algorithm}[H]
\caption{单模数筛表构建 \textsc{BuildTable}}
\label{alg:table}
\begin{algorithmic}[1]
\Require 模数 $m$（奇素数 $p$，或合数 $m = 9$）
\Ensure 周期 $L$ 与布尔表 $G:\ \mathbb{Z}_L \to \{0,1\}$
\State $o \gets 1$;\quad $c \gets 2 \bmod m$
\While{$c \neq 1$}\Comment{求 $o = \Ord_m(2)$，式~\eqref{eq:ord}}
  \State $c \gets c \cdot 2 \bmod m$;\quad $o \gets o + 1$
\EndWhile
\State $L \gets \operatorname{lcm}(m, o)$\Comment{素数时 $L = p\cdot o$，式~\eqref{eq:crt-L}}
\State $\mathit{small}[0] \gets 1 \bmod m$
\For{$i = 0$ \textbf{to} $o-2$}\Comment{幂表 $\mathit{small}[i] = 2^i \bmod m$}
  \State $\mathit{small}[i+1] \gets \mathit{small}[i] \cdot 2 \bmod m$
\EndFor
\State $\mathit{QR}[0..m-1] \gets$ \textbf{false}
\For{$x = 0$ \textbf{to} $m-1$}\Comment{剩余集，式~\eqref{eq:qr-def}；$0$ 必须计入}
  \State $\mathit{QR}[\,x \cdot x \bmod m\,] \gets$ \textbf{true}
\EndFor
\For{$t = 0$ \textbf{to} $L-1$}\Comment{展平双周期，式~\eqref{eq:table-def}}
  \State $G[t] \gets \mathit{QR}\bigl[\,(\mathit{small}[t \bmod o] - t) \bmod m\,\bigr]$
\EndFor
\State \Return $(L,\ G)$
\end{algorithmic}
\end{algorithm}

\paragraph{轮子：把前九个筛烘焙进枚举。}
若先生成全部 $n \equiv 7 \pmod 8$ 再过滤，仍需物化 $C/8$ 个候选。
将模 $8$ 与 $\mathcal{M} = (9,5,7,11,17,13,23,19)$ 的条件合入枚举骨架，
得到周期 $M$ 上的幸存剩余集 $R$ 的递推
\begin{equation}\label{eq:wheel-period}
M_{k+1} = \operatorname{lcm}(M_k, L_{k+1}), \qquad M_0 = 8,\ R_0 = \{7\},
\end{equation}
\begin{equation}\label{eq:wheel-set}
R_{k+1} = \bigl\{\, r + j\,M_k \;:\; r \in R_k,\
0 \le j < M_{k+1}/M_k,\
G_{k+1}\bigl[(r + j M_k) \bmod L_{k+1}\bigr] = 1 \,\bigr\},
\end{equation}
此后扫描只需枚举
\begin{equation}\label{eq:candidate}
n = r + q\,M, \qquad r \in R,\quad q = 0, 1, 2, \dots
\end{equation}
实测 $M = 2{,}677{,}114{,}440$，$|R| = 2{,}177{,}280$，密度
$\rho_{\mathrm{wheel}} = |R|/M = 8.13 \times 10^{-4}$：
约 $99.9\%$ 的候选从未被构造。

\begin{algorithm}[H]
\caption{轮子构建 \textsc{BuildWheel}}
\label{alg:wheel}
\begin{algorithmic}[1]
\Require 模数列表 $\mathcal{M} = (9, 5, 7, 11, 17, 13, 23, 19)$
\Ensure 周期 $M$ 与剩余集 $R \subseteq \mathbb{Z}_M$（式~\eqref{eq:candidate} 的枚举骨架）
\State $M \gets 8$;\quad $R \gets \{7\}$\Comment{式~\eqref{eq:mod8.1}；偶数支路已证无解}
\For{每个 $m \in \mathcal{M}$}
  \State $(L, G) \gets \Call{BuildTable}{m}$
  \State $s \gets L / \gcd(M, L)$\Comment{周期放大倍数，式~\eqref{eq:wheel-period}}
  \State $R' \gets \emptyset$
  \For{每个 $r \in R$}
    \For{$j = 0$ \textbf{to} $s - 1$}
      \If{$G[(r + jM) \bmod L] = 1$}\Comment{保留满足该模数筛的剩余类}
        \State $R' \gets R' \cup \{\,r + jM\,\}$
      \EndIf
    \EndFor
  \EndFor
  \State $R \gets R'$;\quad $M \gets M \cdot s$\Comment{式~\eqref{eq:wheel-set}}
\EndFor
\State \Return $(M,\ R)$
\end{algorithmic}
\end{algorithm}

\begin{algorithm}[H]
\caption{主筛与精确复核 \textsc{Sieve}}
\label{alg:sieve}
\begin{algorithmic}[1]
\Require 上界 $C$；轮子 $(M, R)$（算法~\ref{alg:wheel}）；
\Statex \hspace{\algorithmicindent} 筛表集 $\{(L_i, G_i)\}_{i=1}^{n_f}$
（奇素数 $23 < p \le 421$，共 $n_f = 73$ 张，算法~\ref{alg:table}）
\Ensure $\{\, n < C : 2^n - n \text{ 为完全平方} \,\}$
\State $\mathit{sols} \gets \{1\}$\Comment{$2^1 - 1 = 1 = 1^2$；偶数支路已证无解}
\For{$q = 0$ \textbf{to} $\lceil C/M \rceil - 1$}\Comment{块索引；块间相互独立}
  \State $N \gets \{\, r + qM : r \in R,\ r + qM < C \,\}$\Comment{本块候选，式~\eqref{eq:candidate}}
  \For{$i = 1$ \textbf{to} $n_f$}
    \State $N \gets \{\, n \in N : G_i[\,n \bmod L_i\,] = 1 \,\}$\Comment{取模 + 查表，式~\eqref{eq:equiv}}
    \If{$N = \emptyset$} \textbf{break}\Comment{杀空早退（常态）}
    \EndIf
  \EndFor
  \For{每个 $n \in N$}\Comment{幸存者：期望 $\approx 0$ 个，式~\eqref{eq:expected-survivors}}
    \If{$\bigl\lfloor \sqrt{2^n - n} \bigr\rfloor^{2} = 2^n - n$}\Comment{任意位长精确开方，式~\eqref{eq:exact}}
      \State $\mathit{sols} \gets \mathit{sols} \cup \{n\}$
    \Else\Comment{假阳性：以素数见证复核即否定证书，式~\eqref{eq:euler}}
      \State 记录 $\bigl(p_i,\ a = (2^n - n) \bmod p_i\bigr)$ 供独立复核
    \EndIf
  \EndFor
\EndFor
\State \Return $\mathit{sols}$
\end{algorithmic}
\end{algorithm}

\paragraph{主筛。}
按块 $q$ 扫描，块内候选即式~\eqref{eq:candidate}，逐张过 $n_f$ 张筛表。
数组化实现中第~6 行仅两条向量化指令（取模、花式索引）；由于数组按每筛
$\approx 1/2$ 几何收缩（式~\eqref{eq:count}），总有效过筛次数
$\approx \sum_{k \ge 0} 2^{-k} \approx 2$ 次/候选，故每候选成本为常数个
“取模 + 查表”。块间无依赖、无通信，天然可并行（实现见附录~\ref{app:code}）。

\paragraph{精确复核。}
幸存者用任意位长精确开方判定
\begin{equation}\label{eq:exact}
\bigl\lfloor \sqrt{2^n - n} \bigr\rfloor^{2} \;=\; 2^n - n,
\end{equation}
（Python 标准库 \texttt{math.isqrt}，或 \texttt{gmpy2.is\_square}）。
筛后密度估计为
\begin{equation}\label{eq:density}
\rho \;\approx\;
\underbrace{8.13 \times 10^{-4}}_{\rho_{\mathrm{wheel}}}
\;\cdot\; \prod_{23 < p \le 421} \frac{p+1}{2p}
\;\approx\; 1.5 \times 10^{-25},
\end{equation}
（各筛独立性为启发式，仅影响效率、不影响正确性），故
\begin{equation}\label{eq:expected-survivors}
\mathbb{E}\bigl[\,\#\text{假阳性} \,\big|\, n < C\,\bigr]
\;=\; \rho\, C \;\approx\; 1.5 \times 10^{-10}
\qquad (C = 10^{15}),
\end{equation}
精确复核的执行次数期望为零。实测 $C = 10^{12}$：幸存者恰为 $\{7\}$
（$2^7 - 7 = 121 = 11^2$），零假阳性，与式~\eqref{eq:expected-survivors} 在
$C = 10^{12}$ 处的预言 $1.5\times10^{-13}$ 相符。

\subsection{复杂度与数据漏斗}

\begin{table}[H]
\centering
\caption{$C = 10^{12}$ 的数据漏斗（单核实测 27\,s）}
\begin{tabular}{lrl}
\hline
阶段 & 数量 & 依据 \\
\hline
全体整数 $n < C$ & $10^{12}$ & --- \\
$n \equiv 7 \pmod 8$ & $1.25 \times 10^{11}$ & 式~\eqref{eq:mod8.1} \\
轮子物化候选 & $8.13 \times 10^{8}$ & 式~\eqref{eq:wheel-set} \\
$73$ 筛后幸存 & $1$（$n = 7$） & 式~\eqref{eq:equiv} + 式~\eqref{eq:exact} \\
\hline
\end{tabular}
\end{table}

建表每模 $O(L)$，$73$ 张表合计不足 $2$\,MB，构建不足 $0.5$\,s；轮子一次性
构建不足 $0.5$\,s；主筛候选数 $\rho_{\mathrm{wheel}} C$，每候选常数次取模
加查表，单核实测吞吐 $3.7\times10^{10}$ 个 $n$/s，$C = 10^{15}$ 约需
$7.5$ CPU 小时、内存不足 $100$\,MB；块间独立，多进程近线性加速
（附录~\ref{app:code}）。

\paragraph{正确性小结。}
合法性（式~\eqref{eq:necessity}）、周期性（式~\eqref{eq:ord}--\eqref{eq:equiv}）、
单筛精确对半（式~\eqref{eq:count}）、无假阴性、证书可独立复核
（式~\eqref{eq:euler}）均为定理；唯复合密度（式~\eqref{eq:density}）为
启发式，且它只决定幸存者数量（即精确复核的次数），不影响输出正确性：
若独立性假设完全失真，最坏后果是多复核若干个大数，结论不会错。

\appendix

\section{代码附件：多进程并行实现}\label{app:code}
主文不展开并行细节，仅说明任务划分：任务 = 块区间 $[q_0, q_1)$
（算法~\ref{alg:sieve} 外层循环切片），零通信、只读共享筛表。完整可运行代码如下。

\begin{lstlisting}[language=Python,
  caption={psieve.py: 模筛 + 多进程（任务 = 块区间）},
  label={lst:psieve}]
# -*- coding: utf-8 -*-
# 用法: python psieve.py [C] [worker数]
import math, os, sys, time
import multiprocessing as mp
import numpy as np

def good_table(m):  # 算法 1 (BuildTable)
    o, cur = 1, 2 % m
    while cur != 1:
        cur = cur * 2 % m; o += 1
    L = m * o // math.gcd(m, o)
    small, c = np.empty(o, np.int64), 1 % m
    for i in range(o):
        small[i] = c; c = c * 2 % m
    t = np.arange(L, dtype=np.int64); pw = small[t % o]
    Q = np.zeros(m, bool)
    for x in range(m): Q[x * x % m] = True   # 0 必须计入 QR
    return L, Q[(pw - t) % m]

def build_wheel():  # 算法 2 (BuildWheel)
    M, S = 8, np.array([7], dtype=np.int64)
    for m in (9, 5, 7, 11, 17, 13, 23, 19):
        L, G = good_table(m)
        steps = L // math.gcd(M, L)
        tvec = np.arange(steps, dtype=np.int64) * M
        out, CH = [], max(1, 3_000_000 // steps)
        for i in range(0, S.size, CH):
            cand = S[i:i+CH][:, None] + tvec[None, :]
            out.append(cand[G[cand % L]])
        S = np.concatenate(out); M *= steps
    return M, S

def primes_upto(n):
    s = np.ones(n, bool); s[:2] = False
    for d in range(2, int(n**.5)+1):
        if s[d]: s[d*d::d] = False
    return [int(p) for p in np.nonzero(s)[0]]

def is_square(x):  # 式 (15): 任意位长精确判定
    s = math.isqrt(x); return s * s == x

W_M1 = W_S1 = W_FILT = None
def init_tables():
    global W_M1, W_S1, W_FILT
    if W_M1 is not None: return
    W_M1, W_S1 = build_wheel()
    W_FILT = [good_table(p) for p in primes_upto(430) if p > 23]

def sweep_range(task):  # 算法 3 的一个任务: 块区间 [q0, q1)
    q0, q1, C = task    # imap_unordered 不自动解包, 故收单个元组
    out = []
    for q in range(q0, q1):
        N = W_S1 + q * W_M1            # 式 (12): 本块候选
        for L, G in W_FILT:            # 73 张筛表
            N = N[G[N % L]]            # 式 (9): 取模 + 查表
            if N.size == 0: break
        if N.size: N = N[N < C]
        if N.size: out.append(N)
    return np.concatenate(out) if out else np.zeros(0, np.int64)

def main():
    C = int(sys.argv[1]) if len(sys.argv) > 1 else 10**15
    workers = (int(sys.argv[2]) if len(sys.argv) > 2
               else len(os.sched_getaffinity(0))
               if hasattr(os, "sched_getaffinity") else os.cpu_count() or 4)
    t0 = time.perf_counter()
    init_tables()
    qN = (C + W_M1 - 1) // W_M1
    CHUNK = 32
    tasks = [(q0, min(q0 + CHUNK, qN), C) for q0 in range(0, qN, CHUNK)]
    # Linux 用 fork (表写时复制共享); macOS/Windows 用 spawn (initializer 重建表)
    if sys.platform.startswith("linux") and "fork" in mp.get_all_start_methods():
        ctx = mp.get_context("fork")
    else:
        ctx = mp.get_context("spawn")
    print(f"C={C:,} | blocks={qN:,} | tasks={len(tasks):,} | "
          f"workers={workers} ({ctx.get_start_method()})", flush=True)
    out, t1 = [], time.perf_counter()
    with ctx.Pool(workers, initializer=init_tables) as pool:
        for arr in pool.imap_unordered(sweep_range, tasks, chunksize=1):
            out.append(arr)
    sv = np.concatenate(out) if out else np.zeros(0, np.int64)
    sols = [1]  # n=1: 2-1=1; 偶数支路已证无解
    for n in map(int, sv):  # 幸存者精确复核 (式 15)
        if is_square((1 << n) - n): sols.append(n)
    print(f"n < {C:,}: 2^n-n 为完全平方 <=> n in {sorted(sols)} "
          f"[sweep {time.perf_counter()-t1:,.0f}s, total "
          f"{time.perf_counter()-t0:,.0f}s]", flush=True)

if __name__ == "__main__":
    main()
\end{lstlisting}

\section{结论}
Ridout 方法只说明解集有限；单位约化加显式 Matveev 定理则给出了
可计算的数值上界，即定理 \ref{thm:bound}；依据平方数模 $m$ 的值在数值上界内筛选，可行解仅有 $n=1,7$。

\begin{thebibliography}{99}

\bibitem{Matveev}
E.~M. Matveev, An explicit lower bound for a homogeneous rational linear form
in logarithms of algebraic numbers II,
\emph{Izvestiya: Mathematics} \textbf{64} (2000), 1217--1269.

\bibitem{NesterenkoWaldschmidt}
Yu.~Nesterenko and M.~Waldschmidt, On the approximation of the values of
exponential function and logarithm by algebraic numbers,
\emph{arXiv preprint} math/0002047 (2000).
\url{https://arxiv.org/abs/math/0002047}

\end{thebibliography}

\end{document}
